\section*{Part III: Boards, Investors, and Founder Control}
\addcontentsline{toc}{section}{Part III: Boards, Investors, and Founder Control}

\subsection{14. Control has several instruments}

Economic ownership percentage, shareholder vote thresholds, board appointment rights, reserved-matter consents, liquidation preferences, information rights, transfer restrictions and debt covenants can all affect control. A shareholder may have a minority stake yet block issuing new shares, taking debt, changing the business, hiring a senior leader or selling the company. A founder can retain board influence but face lender covenants preventing distributions or capital projects.

Evaluate the entire rights package. ``Founder still owns a majority'' is not a complete control analysis.

\subsection{15. An illustrative dilution sequence}

\textbf{[Model]} Founder F owns 80 shares and cofounder C owns 20. A new investor contributes \textnormal{₹}4 crore at \textnormal{₹}16 crore pre-money, giving a \textnormal{₹}20-crore post-money valuation and \textbf{20\% post-money ownership}, assuming no option pool or convertibles. F falls to 64\%, C to 16\%. At a subsequent round of \textnormal{₹}10 crore at \textnormal{₹}40 crore pre-money, the new investor receives 20\%; existing stakes multiply by 80\%: F 51.2\%, C 12.8\%, 1st investor 16\%, 2nd investor 20\%.

F remains a narrow economic majority but may already need investor votes for protected decisions. Future option pools and financing can dilute below 50\%. Control must be simulated on a fully diluted basis and by actual contract rights.

\subsection{16. The option-pool shuffle}

An investor may ask for an employee option pool to be created \textbf{before} the financing rather than after, shifting more dilution onto current holders. Compare effective founder ownership and total investor ownership under both structures. Examine who bears the cost of pool top-ups in future rounds, whether unissued pool shares affect fully diluted calculations, and what happens when employees leave.

A high pre-money headline valuation can be economically less attractive after aggressive pool expansion, preference terms and veto rights are considered.

\subsection{17. Board seats and observer rights}

Board rights affect budgets, major transactions, financing, strategy, audit and hiring. An investor director may bring expertise and protect capital but may also have different liquidity preferences. An observer can gain information and influence without the same formal seat; confidentiality and conflicts still matter. Define board size, appointment/removal, quorum, chair, deadlock and committee structure in lawful documentation.

Beware of quorum rules under which 1 absent nominee can repeatedly prevent valid meetings. Include reasonable reconvening mechanisms consistent with applicable requirements.

\subsection{18. Reserved matters are the real negotiation terrain}

Common investor-protected matters may include new securities, changes to share rights, acquisitions, major borrowing, related-party transactions, budgets over specified limits, sale of assets, changes to business, liquidation, founder compensation or departures. Some protections are reasonable investor safeguards; others can turn a minority investor into an operational controller.

Negotiate \textbf{materiality thresholds, duration, consent response times, good-faith standards, ordinary-course carve-outs, cure procedures and sunset rules}. ``Every contract above \textnormal{₹}1 lakh needs investor permission'' is likely unworkable for a growing institution. ``Any undisclosed related-party transaction requires independent approval'' can be protective.

\subsection{19. Shareholder agreements and constitutional consistency}

Investor terms may appear in a term sheet, share-subscription agreement, shareholders' agreement and amended articles. Their enforceability and interplay under Indian law requires counsel; do not assume a private contract automatically overrides statutory rules or constitutional documents. The closing checklist must specify which protections require article amendments, class approvals, board actions, stamp/filing steps and updated registers.

A founder should never sign economic terms and promise to ``sort out governance later.'' Governance is part of price.

\subsection{20. Liquidation preference versus ordinary ownership}

A $1\times$ non-participating preference may allow an investor to take back invested capital or convert to common in a qualifying exit. Participating preferred may receive capital 1st and then share residual proceeds; caps and seniority matter. Multiple rounds can create waterfalls that surprise founders at modest exit values. Anti-dilution terms can shift ownership after down rounds. Redemption rights can create cash pressure even when no buyer exists.

Run low, medium and high exit cases, not just a triumphant unicorn scenario. Compare founder proceeds and company solvency under the real instrument language with counsel.

\subsection{21. Drag-along, tag-along and transfer restrictions}

A drag-along can force consenting holders to sell if defined conditions are met; a tag-along permits participation in specified transfers. Rights of 1st refusal, co-sale, founder lock-in, vesting, good-leaver/bad-leaver and transfer restrictions affect founder liquidity and negotiations. Negotiate trigger thresholds, price protection, representations and personal liability, as well as who decides on the sale.

A founder seeking permanent control should not accept a drag provision that gives a minority investor a unilateral route to force an early sale without carefully understanding it.

\subsection{22. Down rounds and survival financing}

When cash is scarce, founders may accept poor terms not because they are irrational but because failure is near. Protect optionality \textbf{before} the crisis through reserves, spending discipline, recurring revenue, transparent books and several financing alternatives. At a down round, assess preference seniority, anti-dilution, pay-to-play, employee retention, governance reset and whether fresh capital actually funds a viable operating plan.

Control preserved on paper is worthless if the company cannot meet payroll. The sound objective is durable entrepreneurial agency, not an absolute numerical ownership target at any price.

\subsection{23. Information rights and confidentiality}

Investors may request monthly reporting, budget variance, cohort unit economics, bank reconciliations, material customer contracts and board packets. These can improve discipline, but the organization should protect student records, client data, unpublished research, privileged legal materials and competitively sensitive information. Investor ownership does not automatically grant unrestricted access to every personal dataset.

Define authorized reporting, confidentiality, security standards and board access practices, adapting to statutory rights and contracts.

\subsection{24. When to refuse VC or structure it elsewhere}

VC can be appropriate for scalable software, platform infrastructure or research technology with a credible venture exit. It is less naturally aligned with a bounded premium teaching business or a donor-supported academic institution whose cash flows and purpose differ. The founder may preserve the academic core while financing a distinct product company, if legal rights, IP, transfer pricing and related-party governance are clean.

Do not create subsidiaries merely to hide risk or move shareholder value contrary to fiduciary duties. Structure changes require proper approvals and fairness to existing stakeholders.
